\documentclass[border=8pt,varwidth=190mm]{standalone}
% Task 024a -- shared figure style for the Matrix Operator Refactor benchmark figures.
%
% Every figure in this directory is a standalone LaTeX document that \input{}s this
% file and reads its data from tables/*.csv (generated by
% ../../scripts/figures_024a_prepare.jl).  Build with ../../scripts/figures_024a_build.sh.
%
% Palette: four categorical slots, validated colourblind-safe on a white surface
% (all-pairs mode: lightness band PASS, chroma floor PASS, worst CVD dE 8.4,
% worst normal-vision dE 16.3, all four >= 3:1 contrast).  Series are additionally
% separated by mark shape, so identity never rests on colour alone.
% Slots are assigned in fixed order and never cycled or recoloured by rank.

\usepackage[T1]{fontenc}
\usepackage[scaled]{helvet}
\renewcommand{\familydefault}{\sfdefault}
\usepackage{amsmath}
\usepackage{xcolor}
\usepackage{pgfplots}
\usepackage{pgfplotstable}
\usepgfplotslibrary{groupplots}
\usetikzlibrary{patterns}
\pgfplotsset{compat=1.18}

% ---- categorical slots (fixed order) -------------------------------------
\definecolor{fmSone}{HTML}{2A78D6}   % blue
\definecolor{fmStwo}{HTML}{EB6834}   % orange
\definecolor{fmSthree}{HTML}{199E70} % aqua
\definecolor{fmSfour}{HTML}{4A3AA7}  % violet
% ---- chrome and ink ------------------------------------------------------
\definecolor{fmInk}{HTML}{0B0B0B}
\definecolor{fmInkTwo}{HTML}{52514E}
\definecolor{fmMuted}{HTML}{898781}
\definecolor{fmGrid}{HTML}{E1E0D9}
\definecolor{fmAxis}{HTML}{C3C2B7}
% ---- reserved status colour: gates, rooflines, break-even rules ----------
\definecolor{fmCritical}{HTML}{D03B3B}

% Same reference-rule look as the \pgfplotsset fmrule style below, but usable from
% \draw inside an axis (needed on ybar axes, where an \addplot would become a bar).
\tikzset{fmruleline/.style={color=fmCritical, line width=0.7pt, dash pattern=on 3pt off 2pt}}

\newcommand{\fmtitle}[1]{{\bfseries\small\color{fmInk}#1}\par\vspace{2pt}}
\newcommand{\fmnote}[1]{\vspace{1pt}{\scriptsize\color{fmInkTwo}#1}\par}

\pgfplotsset{
  % recessive grid and axes, thin marks, no chart junk
  fmbase/.style={
    width=72mm, height=52mm,
    scale only axis,
    unbounded coords=jump,          % 'nan' in the tables is skipped, never drawn as 0
    every axis plot/.append style={line width=0.9pt, mark size=2.1pt},
    tick label style={font=\scriptsize, color=fmInkTwo},
    label style={font=\scriptsize, color=fmInkTwo},
    title style={font=\small\bfseries, color=fmInk, yshift=1pt},
    legend style={font=\scriptsize, draw=fmAxis, fill=white, fill opacity=0.92,
                  text opacity=1, inner sep=2pt, row sep=-1pt, /tikz/every even column/.append style={column sep=4pt}},
    legend cell align=left,
    grid=major,
    grid style={fmGrid, line width=0.35pt},
    axis line style={fmAxis, line width=0.5pt},
    tick style={fmAxis, line width=0.5pt},
    major tick length=2pt,
    minor tick length=1pt,
    log ticks with fixed point,
  },
  % series slots: colour + mark together (composite encoding)
  fmA/.style={color=fmSone,   mark=*},
  fmB/.style={color=fmStwo,   mark=square*},
  fmC/.style={color=fmSthree, mark=triangle*},
  fmD/.style={color=fmSfour,  mark=diamond*},
  % second channel of the same entity: same colour + mark, dashed, hollow
  fmalt/.style={dashed, mark options={fill=white}},
  % reference rules (break-even, roofline, memory gate) -- never a data series
  fmrule/.style={color=fmCritical, line width=0.7pt, dash pattern=on 3pt off 2pt,
                 mark=none, forget plot},
  % grouped bars: fmbaraxis on the axis (positions the groups, sets the width via
  % the tikz key), fmbars on each \addplot (flat fill, no outline, no marks).
  % 'bar width' lives in /tikz, so it must be spelled out at axis level.
  % log origin y=infty anchors log-scale bars to the axis floor; the pgfplots
  % default (10^0) would draw them from y=1 and misrepresent sub-1 values.
  fmbaraxis/.style={ybar, /tikz/bar width=3.4pt, enlarge x limits=0.14,
                    log origin y=infty},
  fmbars/.style={draw=none, mark=none},
}

\begin{document}

\fmtitle{Fig.\ 3 --- The 019 tuning gains stage by stage, what is left between M2L and the roofline, and what the 023 step costs}

\begin{tikzpicture}
\begin{groupplot}[
  fmbase,
  width=47mm, height=46mm,
  group style={group size=2 by 2, horizontal sep=16mm, vertical sep=12mm},
  ymode=log,
  fmbaraxis,
]

% --- (a) per-stage device time --------------------------------------------
\nextgroupplot[title={(a) device stage time}, ylabel={time (ms)},
               xtick={1,2,3,4,5}, xticklabels={B2M,M2M,M2L,L2L,L2B},
               ymin=0.5, ymax=6000, legend pos=north west, legend columns=1]
\addplot[fmbars, fill=fmSone] table[x=stage_index, y=baseline_022_ms, col sep=comma] {tables/fig03a_stages.csv};
\addlegendentry{022 baseline}
\addplot[fmbars, fill=fmStwo] table[x=stage_index, y=post_abd_ms, col sep=comma] {tables/fig03a_stages.csv};
\addlegendentry{after 019 A+B+D}
\addplot[fmbars, fill=fmSthree] table[x=stage_index, y=post_e_ms, col sep=comma] {tables/fig03a_stages.csv};
\addlegendentry{after 019 E (final)}
\draw[fmruleline] (axis cs:0.4,10) -- (axis cs:5.6,10);
\node[anchor=south east, font=\tiny, color=fmCritical] at (rel axis cs:0.99,0.35)
  {$\approx$10\,ms bandwidth roofline};

% --- (b) 019 evaluation beside one-shot setup; not the 023 recurring loop ---
\nextgroupplot[title={(b) 019 execution and setup costs}, ylabel={time (s)},
               xtick={1,2,3}, xticklabels={device\\exec, host list\\build, state\\build},
               xticklabel style={align=center, font=\tiny},
               ymin=0.05, ymax=12]
\addplot[fmbars, fill=fmSone] table[x=phase_index, y=baseline_022_s, col sep=comma] {tables/fig03b_lifecycle.csv};
\addplot[fmbars, fill=fmStwo] table[x=phase_index, y=post_abd_s, col sep=comma] {tables/fig03b_lifecycle.csv};
\addplot[fmbars, fill=fmSthree] table[x=phase_index, y=post_e_s, col sep=comma] {tables/fig03b_lifecycle.csv};

% --- (c) 023 one-shot setup vs the recurring step it enables ---------------
\nextgroupplot[title={(c) 023 one-shot setup vs recurring step}, ylabel={time (s)},
               xtick={1,2}, xticklabels={$n=10^4$, $n=10^5$},
               xticklabel style={font=\tiny}, ymin=0.01, ymax=800,
               legend pos=north west, legend columns=2]
\addplot[fmbars, fill=fmSone] table[x=n_index, y=cpu_construct_s, col sep=comma] {tables/fig03d_oneshot_vs_recurring.csv};
\addlegendentry{CPU construct}
\addplot[fmbars, fill=fmStwo] table[x=n_index, y=cpu_step_s, col sep=comma] {tables/fig03d_oneshot_vs_recurring.csv};
\addlegendentry{CPU step}
\addplot[fmbars, fill=fmSthree] table[x=n_index, y=gpu_construct_s, col sep=comma] {tables/fig03d_oneshot_vs_recurring.csv};
\addlegendentry{GPU construct}
\addplot[fmbars, fill=fmSfour] table[x=n_index, y=gpu_step_s, col sep=comma] {tables/fig03d_oneshot_vs_recurring.csv};
\addlegendentry{GPU step}
\node[anchor=west, font=\tiny, color=fmCritical, align=left]
  at (axis cs:1.09,23.7) {23.7\,s: CUDA JIT,\\not setup cost};

% --- (d) 023 resident step split ------------------------------------------
\nextgroupplot[title={(d) 023 resident step split}, ylabel={time (s)},
               xtick={1,2,3}, xticklabels={update, lifecycle, finalize},
               xticklabel style={font=\tiny}, ymin=5e-4, ymax=0.5,
               legend pos=north west, legend columns=1]
\addplot[fmbars, fill=fmSone] table[x=phase_index, y=n10000_s, col sep=comma] {tables/fig03c_step_split.csv};
\addlegendentry{$n=10^4$}
\addplot[fmbars, fill=fmStwo] table[x=phase_index, y=n100000_s, col sep=comma] {tables/fig03c_step_split.csv};
\addlegendentry{$n=10^5$}

\end{groupplot}
\end{tikzpicture}

\fmnote{\textbf{Machine / regime:} \texttt{m13h-1-1}, NVIDIA H200, CUDA 12.8, Julia 1.11.7. Panels (a,b):
$n=10^5$, $\ell=4$, $P=4$, \texttt{ConstantPAnalyticStencil(4,1e-4)}, concat chunk $2^{17}$.
Panels (c,d): task 023 production-integration radix path, $P=4$, $\ell=4$, Julia \texttt{threads=8}.
\textbf{Data:} \texttt{operator\_performance\_tuning/cuda\_022\_baseline.csv},
\texttt{m13h-1-1/cuda\_019\_phase\{ABD,E\}\_throughput.csv},
\texttt{production\_integration/benchmark\_023\_m13h-1-1\_20260714-233159.csv}
(tables \texttt{fig03a\_*}, \texttt{fig03b\_*}, \texttt{fig03d\_*} for panel (c),
\texttt{fig03c\_*} for panel (d)).
\textbf{Reads as:} M2L fell $1698\to66$\,ms ($\approx25\times$) and M2M/L2L $\approx150\to8$\,ms
($\approx19\times$), taking the whole device lifecycle $1.863\to0.116$\,s ($\approx16\times$). M2L is now
$\approx6\times$ above the memory-bandwidth roofline, down from $\approx100\times$. The largest remaining
device stage is L2B at 50\,ms, which the tuning did not touch --- and which is not a far-field operator at
all: L2B evaluates the \emph{nearfield} direct pairs as well as the local expansions, so most of that bar is
$O(N^2)$-within-the-complement work that no operator strategy can remove. In the pre-023
reconstruct-every-call workflow, the 0.35\,s host list build and 0.62\,s state build would have dominated
the 0.116\,s device evaluation. Panel (d) shows the actual 023 recurring step: device-side update replaces
those builds, so neither setup bar belongs in the recurring loop. Panel (c) prices the remaining one-shot
cost against the step it enables: on the GPU the representative cache construction is 0.126\,s
($n=10^5$) against a 0.086\,s step, so device residency pays for itself within about two steps, while on
the CPU an 0.41\,s construction sits in front of a 24.6\,s step that is not worth taking at all
(Fig.\ 8).
\textbf{Caveats:} the 022 ``baseline'' stage times are \emph{single-shot} while both post-tuning columns
are min-of-3, so the $\approx25\times$ M2L figure mixes timing methods. The post-E B2M bar (14\,ms vs
1.4\,ms before) is measurement noise, not a regression --- task 019 recorded B2M jitter of 8--26\,ms across
sweep rows, and the other two chunk widths in the same phase-E sweep read 1.4\,ms. Panel (c)'s GPU
construction bar at $n=10^4$ (23.7\,s) is \emph{not} a setup cost: that case runs first in the 023 script, so
it absorbs one-time CUDA kernel compilation --- the same host reads 0.126\,s at $n=10^5$ and
\texttt{m13h-2-1} reads 24.7\,s then 0.030\,s for the same pair. Read the $n=10^5$ bar as the
representative construction cost. \textbf{Scope:} panel
(b)'s device exec and the step bars of panels (c,d) are complete \texttt{fmm!} evaluations --- far field \emph{and}
nearfield, the latter inside L2B --- so they are end-to-end influence, not a far-field-only subset; task 024
validates the same pipeline against an all-bodies \texttt{direct!} reference. The chunk width itself
changes the lifecycle minimum by at most $\approx6\%$; $2^{17}$ was kept for least scratch, not because it
was fastest.}

\end{document}
